% AA-39 | Audyt I: samodzielny dowód wymiaru K[U,V], doprecyzowanie łańcuchów i odsyłaczy.
\chapter{Wymiar Krulla i łańcuchy ideałów pierwszych}
\label{aa:ch:krull-dimension}

\section{Dlaczego mierzymy długości łańcuchów?}
W przestrzeni afinicznej nad ciałem $K$ prosta ma jeden swobodny
parametr, a płaszczyzna dwa. Dla pierścieni zamiast współrzędnych
można badać ciągi coraz większych ideałów pierwszych. Pozwala to
zdefiniować wymiar także dla pierścieni, które nie są pierścieniami
wielomianów, a nawet mają dzielniki zera. W całym rozdziale pierścienie
są przemienne, mają jedynkę i są niezerowe, o ile wyraźnie nie
postanowimy inaczej. Wprowadzenie ideałów pierwszych i maksymalnych
znajduje się w Rozdziale~\ref{aa:ch:pierscienie}.

\begin{definicja}[Łańcuch ideałów i wymiar Krulla]\label{aa:def:krull-dimension}
Niech $A$ będzie niezerowym pierścieniem przemiennym z jedynką.
\emph{Łańcuchem ideałów pierwszych długości $r$} nazywamy ciąg
\[
 \mathfrak p_0\subsetneq\mathfrak p_1\subsetneq\cdots
 \subsetneq\mathfrak p_r
\]
ideałów pierwszych pierścienia $A$, gdzie $r\in\mathbb N_0$.
\emph{Wymiar Krulla} $\dim A$ jest kresem górnym zbioru długości
wszystkich takich łańcuchów; może mieć wartość $+\infty$.
Dla ideału pierwszego $\mathfrak p$ definiujemy jego \emph{wysokość}
$\operatorname{ht}(\mathfrak p)$ jako kres górny długości łańcuchów
ideałów pierwszych kończących się w $\mathfrak p$.
\end{definicja}

Długość oznacza liczbę ścisłych zawierań, a nie liczbę występujących
ideałów. W szczególności pojedynczy ideał pierwszy tworzy łańcuch
długości $0$. Istnienie ideału maksymalnego w każdym niezerowym
pierścieniu przemiennym z jedynką wynika z lematu Zorna (w ogólności
wykorzystujemy tu aksjomat wyboru), więc $\dim A\ge0$.
Nie wprowadzamy tu wymiaru pierścienia zerowego; wyłączenie go
zapobiega niejasnej konwencji dla pustego spektrum.

Jeżeli $\mathfrak p$ jest ideałem pierwszym, to nie wolno
utożsamiać $\operatorname{ht}(\mathfrak p)$ z liczbą
ideałów pierwszych zawartych w $\mathfrak p$:
wysokość liczy długości \emph{ściśle rosnących łańcuchów}.
Na przykład w $\mathbb Z$ mamy
$\operatorname{ht}((0))=0$ oraz
$\operatorname{ht}((p))=1$ dla każdej liczby pierwszej $p$.
Dla dowolnego pierścienia możliwa jest również wysokość nieskończona.



\begin{wniosek}[Charakterystyka wymiaru zero]
\label{aa:cor:krull-dim-zero}
Niech $A$ będzie niezerowym pierścieniem przemiennym z jedynką.
Wówczas $\dim A=0$ wtedy i tylko wtedy, gdy każdy ideał pierwszy
pierścienia $A$ jest maksymalny.
\end{wniosek}
\begin{proof}
Jeśli $\dim A=0$, nie istnieją dwa różne ideały pierwsze
$\mathfrak p\subsetneq\mathfrak q$. Każdy ideał pierwszy
$\mathfrak p$ zawiera się w pewnym ideale maksymalnym
$\mathfrak m$, który również jest pierwszy; stąd
$\mathfrak p=\mathfrak m$. Odwrotnie, jeśli wszystkie ideały
pierwsze są maksymalne, żadne dwa różne takie ideały nie mogą
pozostawać w relacji ścisłego zawierania. Ponieważ $A$ posiada
co najmniej jeden ideał pierwszy, największa możliwa długość
łańcucha wynosi $0$.
\end{proof}
\begin{przyklad}[Ciało i dziedzina liczb całkowitych]\label{aa:prz:krull-field-integers}
W ciele $K$ jedynym ideałem właściwym jest $(0)$, który jest
pierwszy. Nie można utworzyć łańcucha długości $1$, dlatego
$\dim K=0$.
W $\mathbb Z$ ideał $(0)$ jest pierwszy, ponieważ $\mathbb Z$
nie ma dzielników zera. Każdy niezerowy ideał pierwszy ma postać
$(p)$ dla liczby pierwszej $p$ i jest maksymalny, gdyż
$\mathbb Z/(p)$ jest ciałem. Stąd wszystkie nietrywialne łańcuchy
mają postać $(0)\subsetneq(p)$ i $\dim\mathbb Z=1$.
\end{przyklad}

\section{Ilorazy i lokalizacje}
\begin{twierdzenie}[Ideały pierwsze w ilorazie]\label{aa:tw:krull-quotient-primes}
Niech $A$ będzie pierścieniem przemiennym z jedynką, a $I\subsetneq A$
ideałem. Odwzorowanie $\mathfrak p\mapsto\mathfrak p/I$ jest
zachowującą inkluzję bijekcją między ideałami pierwszymi $A$
zawierającymi $I$ a ideałami pierwszymi $A/I$.
W szczególności $\dim(A/I)$ jest kresem długości łańcuchów
ideałów pierwszych $A$ zawierających $I$.
\end{twierdzenie}
\begin{proof}
Niech $\pi:A\to A/I$ będzie homomorfizmem ilorazowym.
Jeśli $\mathfrak p$ zawiera $I$ i jest pierwszy, to
$(A/I)/(\mathfrak p/I)\cong A/\mathfrak p$ jest dziedziną
całkowitości, a więc $\mathfrak p/I$ jest pierwszy.
Jeśli $\mathfrak q$ jest pierwszy w $A/I$, jego przeciwobraz
$\pi^{-1}(\mathfrak q)$ jest pierwszy w $A$: z
$ab\in\pi^{-1}(\mathfrak q)$ wynika
$\pi(a)\pi(b)\in\mathfrak q$, zatem $\pi(a)\in\mathfrak q$
lub $\pi(b)\in\mathfrak q$. Przeciwobraz zawiera $I$.
Sprawdźmy odwrotność obu konstrukcji. Dla ideału
$\mathfrak p\supseteq I$ równość
$\pi^{-1}(\mathfrak p/I)=\mathfrak p$ wynika z tego, że
$\pi(a)\in\mathfrak p/I$ dokładnie wtedy, gdy
$a\in\mathfrak p$. Dla ideału $\mathfrak q\subseteq A/I$
mamy $\pi(\pi^{-1}(\mathfrak q))=\mathfrak q$, ponieważ
$\pi$ jest surjekcją. Jeśli $\mathfrak p_1\subsetneq
\mathfrak p_2$ zawierają $I$, to ich obrazy pozostają różne:
wybierzmy $a\in\mathfrak p_2\setminus\mathfrak p_1$;
wtedy $a+I$ należy do $\mathfrak p_2/I$, lecz nie należy
do $\mathfrak p_1/I$. Analogicznie działamy w kierunku
przeciwnym przez przeciwobrazy. Tym samym bijekcja zachowuje
i odzwierciedla ścisłe inkluzje, więc przenosi łańcuchy
z zachowaniem długości.
\end{proof}




\begin{wniosek}[Przejście do ilorazu nie zwiększa wymiaru]
\label{aa:cor:krull-quotient-bound}
Niech $A$ będzie niezerowym pierścieniem przemiennym z jedynką,
a $I\subsetneq A$ ideałem. Wówczas $\dim(A/I)\leq\dim A$.
Ogólniej, jeśli $I\subseteq J\subsetneq A$, to
$\dim(A/J)\leq\dim(A/I)$.
\end{wniosek}
\begin{proof}
Każdy łańcuch ideałów pierwszych $A/I$ odpowiada, zgodnie z
Twierdzeniem~\ref{aa:tw:krull-quotient-primes}, łańcuchowi
ideałów pierwszych $A$ zawierających $I$ i mającemu dokładnie
tę samą długość. W szczególności jego długość nie przekracza
$\dim A$. Dla drugiej nierówności rozpatrujemy
$A/J\cong(A/I)/(J/I)$ i stosujemy pierwszą nierówność
do pierścienia $A/I$, który jest niezerowy, gdyż $I\subsetneq A$.
\end{proof}

\begin{wniosek}[Redukcja nie zmienia wymiaru]\label{aa:cor:krull-nilradical}
Niech $A$ będzie niezerowym pierścieniem przemiennym z jedynką,
a $N=\sqrt{(0)}=\{a\in A:\text{ istnieje }m\ge1\text{ takie, że }a^m=0\}$
jego nilradykałem. Wówczas $A/N$ jest niezerowym pierścieniem
i $\dim(A/N)=\dim A$. Ponadto ideały pierwsze tych dwóch
pierścieni odpowiadają sobie bijektywnie, z zachowaniem wszystkich
ścisłych zawierań.
\end{wniosek}
\begin{proof}
Najpierw $N\ne A$: gdyby $1\in N$, dla pewnego $m\ge1$
mielibyśmy $1^m=0$, a wtedy $A$ byłby pierścieniem zerowym.
Niech $\mathfrak p$ będzie dowolnym ideałem pierwszym $A$.
Dla $a\in N$ istnieje $m\ge1$ takie, że $a^m=0\in\mathfrak p$.
Pierwszość $\mathfrak p$, zastosowana kolejno do iloczynów
$a\cdot a^{m-1}$, $a\cdot a^{m-2}$ itd., daje $a\in\mathfrak p$.
Zatem $N\subseteq\mathfrak p$ dla \emph{każdego} ideału
pierwszego $A$. Twierdzenie~\ref{aa:tw:krull-quotient-primes}
ustanawia więc bijekcję między wszystkimi ideałami pierwszymi
$A$ i $A/N$, zachowującą i odzwierciedlającą ścisłe zawieranie.
Przenosi ona każdy łańcuch na łańcuch tej samej długości
w obie strony, skąd wynika równość wymiarów.
\end{proof}
\begin{twierdzenie}[Ideały pierwsze w lokalizacji]\label{aa:tw:krull-localization-primes}
Niech $A$ będzie pierścieniem przemiennym z jedynką,
a $S\subseteq A$ zbiorem multiplikatywnym zawierającym $1$
i niezawierającym $0$. Ideały pierwsze $S^{-1}A$ odpowiadają
bijektywnie ideałom pierwszym $\mathfrak p\subset A$ takim,
że $\mathfrak p\cap S=\varnothing$, przez
$\mathfrak p\mapsto S^{-1}\mathfrak p$ i przeciwobraz względem
$A\to S^{-1}A$. Obie operacje zachowują i odzwierciedlają inkluzję.
\end{twierdzenie}
\begin{proof}
Ułamki formalne i warunek ich równości są opisane
w Rozdziale~\ref{aa:ch:lokalizacja}.
Dla $\mathfrak p\cap S=\varnothing$ oznaczmy przez
$\overline S$ obraz $S$ w $A/\mathfrak p$.
Homomorfizm
\[
\Phi:S^{-1}A\longrightarrow
\overline S^{-1}(A/\mathfrak p),\qquad
\Phi(a/s)=(a+\mathfrak p)/(s+\mathfrak p)
\]
jest dobrze określony: jeśli $a/s=a'/s'$, to istnieje
$t\in S$ z $t(as'-a's)=0$ w $A$ i ten sam warunek
przechodzi do ilorazu. Jest on surjektywny, ponieważ
każdy ułamek po prawej ma postać obrazu $a/s$.
Ponadto $\Phi(a/s)=0$ dokładnie wtedy, gdy dla pewnego
$t\in S$ zachodzi $ta\in\mathfrak p$.
Pierwszość $\mathfrak p$ oraz $t\notin\mathfrak p$ dają
$a\in\mathfrak p$; w drugą stronę jest to bezpośrednie.
Jądro $\Phi$ jest więc równe $S^{-1}\mathfrak p$.
Ponieważ $A/\mathfrak p$ jest dziedziną całkowitości,
a każdy element $\overline S$ jest niezerowy,
lokalizacja $\overline S^{-1}(A/\mathfrak p)$ również
jest niezerową dziedziną. Z izomorfizmu pierwszego wynika,
że $S^{-1}\mathfrak p$ jest ideałem pierwszym.

Jeśli $\mathfrak q$ jest pierwszy w $S^{-1}A$, jego przeciwobraz
$\mathfrak p$ jest pierwszy w $A$ (własność ta wynika z
zachowania iloczynów przez homomorfizm). Nie może zawierać
elementu $s\in S$, ponieważ $s/1$ jest jednostką lokalizacji,
a ideał właściwy nie zawiera jednostek. Dla $a/s\in S^{-1}A$
mamy $a/s\in\mathfrak q$ wtedy i tylko wtedy, gdy
$a/1=(s/1)(a/s)\in\mathfrak q$, czyli $a\in\mathfrak p$.
Stąd $\mathfrak q=S^{-1}\mathfrak p$; ten sam rachunek pokazuje,
że przeciwobraz $S^{-1}\mathfrak p$ równa się $\mathfrak p$:
jeśli $a/1\in S^{-1}\mathfrak p$, istnieje $s\in S$ z
$sa\in\mathfrak p$, a pierwszość i $s\notin\mathfrak p$
dają $a\in\mathfrak p$. Obie operacje są więc odwrotne.
\end{proof}



\begin{wniosek}[Lokalizacja nie zwiększa wymiaru]
\label{aa:cor:krull-localization-bound}
Niech $A$ będzie niezerowym pierścieniem przemiennym z jedynką,
a $S\subseteq A$ zbiorem multiplikatywnym takim, że
$1\in S$ i $0\notin S$. Wtedy
$\dim(S^{-1}A)\le\dim A$.
\end{wniosek}
\begin{proof}
Każdy łańcuch ideałów pierwszych pierścienia $S^{-1}A$
przechodzi przez przeciwobrazy względem homomorfizmu
$A\to S^{-1}A$ na łańcuch ideałów pierwszych $A$.
Na mocy Twierdzenia~\ref{aa:tw:krull-localization-primes}
przeciwobrazy zachowują ścisłość wszystkich zawierań.
Długość łańcucha w lokalizacji nie przekracza zatem
kresu długości łańcuchów w $A$. Przechodząc do kresu
górnego po wszystkich łańcuchach w lokalizacji, otrzymujemy tezę.
\end{proof}
\begin{wniosek}[Wysokość a pierścień lokalny]\label{aa:cor:krull-height-local}
Niech $\mathfrak p$ będzie ideałem pierwszym pierścienia
przemiennego $A$ z jedynką. Wtedy
$\operatorname{ht}(\mathfrak p)=\dim A_{\mathfrak p}$,
gdzie $A_{\mathfrak p}=(A\setminus\mathfrak p)^{-1}A$.
\end{wniosek}
\begin{proof}
Ideały pierwsze $A_{\mathfrak p}$ odpowiadają, na mocy
Twierdzenia~\ref{aa:tw:krull-localization-primes}, dokładnie tym
ideałom pierwszym $\mathfrak q$ pierścienia $A$, które są rozłączne
z $A\setminus\mathfrak p$, czyli spełniają
$\mathfrak q\subseteq\mathfrak p$.
Każdy łańcuch takich ideałów można przedłużyć do
$\mathfrak p$, jeżeli jeszcze na nim nie kończy się.
Stąd kres długości wszystkich takich łańcuchów jest
równy kresowi długości łańcuchów kończących się w $\mathfrak p$.
\end{proof}

\section{Pierwsze obliczenia wymiaru}
\begin{twierdzenie}[Wymiar dziedziny ideałów głównych]\label{aa:tw:krull-pid}
Niech $A$ będzie dziedziną ideałów głównych. Jeżeli $A$ jest
ciałem, to $\dim A=0$; w przeciwnym razie $\dim A=1$.
\end{twierdzenie}
\begin{proof}
W dziedzinie $A$ ideał $(0)$ jest pierwszy. Jeśli $A$ jest
ciałem, innych ideałów pierwszych nie ma.
Jeśli $A$ nie jest ciałem, wybierzmy niezerowy element
nieodwracalny $a\in A$. Właściwy ideał $(a)$ jest zawarty
w pewnym ideale maksymalnym $\mathfrak m$ (lemat Zorna).
Jest on niezerowy, a jako maksymalny również pierwszy.

Pokażemy, że każdy niezerowy ideał pierwszy jest maksymalny.
Niech $\mathfrak p=(p)$ będzie niezerowy i pierwszy,
a $(p)\subseteq (b)\subsetneq A$. Ponieważ $p\in(b)$,
możemy napisać $p=bc$. Pierwszość ideału $(p)$ oznacza,
że $b\in(p)$ lub $c\in(p)$. Jeżeli $b=pu$, to
$p=puc$, a skracanie przez $p\ne0$ daje $1=uc$.
Wtedy $c$ jest jednostką i $(b)=(p)$. Jeżeli zaś
$c=pv$, to podobnie $1=bv$, zatem $(b)=A$,
co jest sprzeczne z założeniem. Nie istnieje więc
właściwy ideał ściśle zawierający $\mathfrak p$.
Każdy niezerowy ideał pierwszy jest maksymalny, więc łańcuchy
mają długość co najwyżej $1$, a $(0)\subsetneq(p)$
pokazuje, że długość $1$ występuje.
\end{proof}

\begin{lemat}[Wymiar pierścienia wielomianów w dwóch zmiennych]
\label{aa:lem:krull-plane-dimension}
Niech $K$ będzie dowolnym ciałem. Wówczas
$\dim K[U,V]=2$. Ponadto każdy niezerowy ideał pierwszy
$\mathfrak p$ pierścienia $K[U,V]$ jest maksymalny albo ma
postać $(f)$ dla pewnego nierozkładalnego $f\in K[U,V]$.
\end{lemat}
\begin{proof}
Oznaczmy $R=K[U,V]$. Ponieważ $K[U]$ jest pierścieniem
z jednoznacznością rozkładu, z lematu Gaussa wynika, że
$R=K[U][V]$ także jest takim pierścieniem. W szczególności
każdy niezerowy element nieodwracalny rozkłada się na
iloczyn elementów nierozkładalnych.

Wykażemy najpierw, że dla niezerowych, nierozkładalnych
i niestowarzyszonych $f,g\in R$ iloraz $R/(f,g)$ jest
skończenie wymiarową przestrzenią liniową nad $K$.
W pierścieniu $K(U)[V]$ wielomiany $f,g$ są względnie
pierwsze. Istotnie, każdy wspólny dzielnik dodatniego stopnia
względem $V$ dawałby na mocy lematu Gaussa wspólny czynnik
w $K[U,V]$; czynnik stopnia zero jest zaś jednostką
w $K(U)[V]$. Z tożsamości Bézouta w pierścieniu
$K(U)[V]$ otrzymujemy wielomiany $a,b\in K(U)[V]$ takie,
że $af+bg=1$. Po usunięciu mianowników dostajemy
niezerowy $h(U)\in K[U]$ należący do $(f,g)$.
Zamieniając role $U$ i $V$, uzyskujemy analogicznie
niezerowy $k(V)\in K[V]$ należący do $(f,g)$.
W ilorazie $R/(f,g)$ można obniżać stopnie osobno względem
$U$ i $V$ za pomocą równań $h(U)=0$ oraz $k(V)=0$:
współczynniki wiodące obu wielomianów są niezerowymi
skalarami z $K$. Zatem klasy jednomianów
$U^iV^j$ dla $0\le i<\deg h$ i $0\le j<\deg k$
rozpinają $R/(f,g)$ nad $K$. Jeżeli któryś z wielomianów
$h,k$ jest stałą niezerową, iloraz jest zerowy, więc
wniosek o skończonym wymiarze pozostaje prawdziwy.

Niech teraz $\mathfrak p\ne(0)$ będzie ideałem pierwszym $R$.
Zawiera on niezerowy element $a$. Rozkładając $a$ na
czynniki nierozkładalne i korzystając z pierwszości
$\mathfrak p$, znajdujemy nierozkładalne $f\in\mathfrak p$.
Jeżeli $\mathfrak p\ne(f)$, istnieje
$g\in\mathfrak p\setminus(f)$. Po rozłożeniu $g$ na
czynniki nierozkładalne przynajmniej jeden z nich,
powiedzmy $g_0$, należy do $\mathfrak p$ i nie jest
stowarzyszony z $f$. Wykazane wyżej rozumowanie daje
$\dim_K R/(f,g_0)<\infty$. Tym bardziej
$R/\mathfrak p$, będący jego ilorazem, jest skończenie
wymiarową $K$-algebrą. Ponieważ $\mathfrak p$ jest
pierwszy, $R/\mathfrak p$ jest dziedziną. Każda niezerowa
skończenie wymiarowa dziedzina nad ciałem jest ciałem:
mnożenie przez dowolny element niezerowy jest injektywnym,
a zatem surjektywnym, endomorfizmem przestrzeni liniowej.
Wobec tego $\mathfrak p$ jest maksymalny.

Każdy niemaksymalny ideał pierwszy $R$ jest więc równy
$(0)$ albo ma postać $(f)$ dla pewnego nierozkładalnego $f$.
Ideał $(f)$ jest pierwszy, ponieważ w pierścieniu z
jednoznacznością rozkładu elementy nierozkładalne są pierwsze.
Jeżeli $(f)\subsetneq\mathfrak q$ dla ideału pierwszego
$\mathfrak q$, to $\mathfrak q$ nie może być innym głównym
ideałem $(g)$ generowanym przez element nierozkładalny:
z $f\in(g)$ wynikałoby, że $f$ i $g$ są stowarzyszone,
a wtedy $(f)=(g)$. Zatem $\mathfrak q$ jest maksymalny.
Wynika stąd, że długość łańcucha ideałów pierwszych
$R$ nie przekracza $2$. Z drugiej strony łańcuch
$(0)\subsetneq(U)\subsetneq(U,V)$ ma długość $2$;
wszystkie trzy ideały są pierwsze na mocy odpowiednich
ilorazów $R$, $K[V]$ i $K$. Zatem $\dim K[U,V]=2$.
\end{proof}

\begin{przyklad}[Krzyżujące się proste]\label{aa:prz:krull-crossed-lines}
Niech $K$ będzie ciałem, $A=K[X,Y]/(XY)$, a $x,y$ oznaczają
klasy zmiennych. W pierścieniu $A$ mamy $xy=0$, choć
$x\ne0$ i $y\ne0$. Jeżeli $\mathfrak p$ jest ideałem
pierwszym $A$, z $xy\in\mathfrak p$ wynika
$x\in\mathfrak p$ lub $y\in\mathfrak p$.
W pierwszym przypadku $\mathfrak p$ pochodzi od ideału
pierwszego w $A/(x)\cong K[Y]$; w drugim od ideału
pierwszego w $A/(y)\cong K[X]$. Rozważmy teraz
\emph{dowolny} łańcuch $\mathfrak p_0\subsetneq\cdots
\subsetneq\mathfrak p_r$. Ponieważ $xy=0\in\mathfrak p_0$,
pierwszość $\mathfrak p_0$ daje $x\in\mathfrak p_0$
lub $y\in\mathfrak p_0$. W pierwszym przypadku każdy
$\mathfrak p_j$ zawiera $x$, w drugim każdy zawiera $y$.
Przez odpowiedniość ideałów pierwszych w ilorazie cały
łańcuch przechodzi więc do łańcucha w $K[Y]$ albo $K[X]$,
którego długość nie przekracza $1$.
Ponieważ $A/(x)\cong K[Y]$ i mamy ścisły łańcuch
$(x)\subsetneq(x,y)$, otrzymujemy $\dim A=1$.
Ideały $(x)$ i $(y)$ są pierwsze, ponieważ odpowiednie
ilorazy są dziedzinami $K[Y]$ i $K[X]$.
Są one minimalne: gdyby ideał pierwszy $\mathfrak q$
był ściśle zawarty w $(x)$, to $y\notin\mathfrak q$
(gdyż $y\notin(x)$), a z $xy=0$ wynikałoby
$x\in\mathfrak q$, zatem $\mathfrak q=(x)$,
sprzeczność. Analogicznie postępujemy z $(y)$.
Dwa składniki geometryczne, opisane równaniami $x=0$
i $y=0$, spotykają się w jednym punkcie, lecz nie zwiększają
wymiaru całości do dwóch. Rysunek~\ref{aa:fig:crossed-lines}
jest schematem dla $K=\mathbb R$; rachunek ideałowy
pozostaje prawdziwy nad dowolnym ciałem.
\end{przyklad}
\begin{figure}[htbp]
\centering
\setlength{\unitlength}{1mm}
\begin{picture}(82,43)
\put(11,19){\vector(1,0){60}}
\put(40,3){\vector(0,1){37}}
\put(73,16){$x$}
\put(43,37){$y$}
\put(54,13){$y=0$}
\put(44,31){$x=0$}
\put(37,15){$0$}
\end{picture}
\caption{Krzywa afiniczna $xy=0$ składa się z dwóch prostych
jednowymiarowych przecinających się w punkcie początkowym.
Nie jest całą dwuwymiarową płaszczyzną afiniczną; rysunek
przedstawia punkty rzeczywiste, a obliczenie wymiaru działa
nad dowolnym ciałem.}
\label{aa:fig:crossed-lines}
\end{figure}


\begin{twierdzenie}[Wymiar iloczynu dwóch pierścieni]
\label{aa:tw:krull-product-dimension}
Niech $A$ i $B$ będą niezerowymi pierścieniami przemiennymi
z jedynką. Wówczas
\[
 \dim(A\times B)=\max\{\dim A,\dim B\}.
\]
Każdy ideał pierwszy $A\times B$ ma dokładnie jedną
z postaci $\mathfrak p\times B$ lub $A\times\mathfrak q$,
gdzie $\mathfrak p$ jest ideałem pierwszym $A$, a
$\mathfrak q$ ideałem pierwszym $B$.
\end{twierdzenie}
\begin{proof}
W pierścieniu $A\times B$ rozważmy elementy
$e=(1_A,0_B)$ i $f=(0_A,1_B)$. Zachodzą równości
$ef=0$ i $e+f=1$. Jeżeli $\mathfrak r$ jest ideałem
pierwszym $A\times B$, to z $ef\in\mathfrak r$ wynika
$e\in\mathfrak r$ lub $f\in\mathfrak r$. Oba nie mogą
należeć do $\mathfrak r$, ponieważ wtedy należałaby doń
jedynka, a ideał pierwszy jest właściwy.

Jeżeli $f\in\mathfrak r$, to $\{0\}\times B\subseteq
\mathfrak r$. Istnieje zatem ideał $\mathfrak p\subseteq A$
taki, że $\mathfrak r=\mathfrak p\times B$; istotnie,
wystarczy wziąć $\mathfrak p=\{a\in A:(a,0)\in\mathfrak r\}$.
Iloraz $(A\times B)/(\mathfrak p\times B)$ jest izomorficzny
z $A/\mathfrak p$, a więc $\mathfrak r$ jest pierwszy
dokładnie wtedy, gdy $\mathfrak p$ jest pierwszy.
Analogicznie przypadek $e\in\mathfrak r$ daje
$\mathfrak r=A\times\mathfrak q$ dla ideału pierwszego
$\mathfrak q$ pierścienia $B$.

Ideały pierwsze należące do różnych opisanych rodzin
nie mogą się zawierać jeden w drugim. Istotnie,
$\mathfrak p\times B$ zawiera $f$, natomiast
$A\times\mathfrak q$ go nie zawiera (gdyż
$1_B\notin\mathfrak q$); w odwrotną stronę stosujemy $e$.
W obrębie każdej rodziny inkluzje ideałów odpowiadają
dokładnie inkluzjom w odpowiednim czynniku, np.
$\mathfrak p_1\times B\subsetneq\mathfrak p_2\times B$
wtedy i tylko wtedy, gdy
$\mathfrak p_1\subsetneq\mathfrak p_2$.
Stąd każdy łańcuch ideałów pierwszych iloczynu leży w jednej
rodzinie, a łańcuchy obu czynników przenoszą się do iloczynu
bez zmiany długości. Biorąc kresy długości, otrzymujemy
wskazaną równość, również gdy któryś wymiar jest nieskończony.
\end{proof}

\begin{przyklad}[Maksymalne ideały o różnych wysokościach]
\label{aa:prz:krull-maximal-varying-heights}
Niech $K$ będzie ciałem,
$A=K[T]$, $B=K[U,V]$, a $R=A\times B$.
Z poprzedniego twierdzenia, obliczenia wymiaru $K[T]$
(Twierdzenie~\ref{aa:tw:krull-pid}) oraz
Lematu~\ref{aa:lem:krull-plane-dimension} wynika,
że $\dim R=\max\{1,2\}=2$.
Rozważmy ideały
$\mathfrak m_1=(T)\times B$ oraz
$\mathfrak m_2=A\times(U,V)$. Są maksymalne, ponieważ
$R/\mathfrak m_i\cong K$ dla $i=1,2$.
Każdy łańcuch ideałów pierwszych kończący się w
$\mathfrak m_1$ ma postać
$\mathfrak p_0\times B\subsetneq\cdots\subsetneq(T)\times B$;
odpowiada więc łańcuchowi w $K[T]$.
Łańcuch $(0)\times B\subsetneq(T)\times B$ istnieje,
a dłuższy nie może istnieć, stąd
$\operatorname{ht}(\mathfrak m_1)=1$.
Podobnie łańcuch
$A\times(0)\subsetneq A\times(U)\subsetneq A\times(U,V)$
pokazuje $\operatorname{ht}(\mathfrak m_2)\ge2$,
a obliczenie wymiaru $K[U,V]$ daje nierówność przeciwną.
Zatem $\operatorname{ht}(\mathfrak m_2)=2$.
W szczególności ideał maksymalny nie musi mieć wysokości
równej wymiarowi całego pierścienia. Zauważmy, że oba
występujące czynniki są dziedzinami, ale ich iloczyn już nie.
\end{przyklad}

\begin{przyklad}[Iloraz może zmniejszyć wymiar]
Niech $A=\mathbb Z$, $I=(6)$. Ideały pierwsze $A/I$
odpowiadają ideałom pierwszym $\mathbb Z$ zawierającym $(6)$.
Wyjaśnijmy, skąd pochodzi pełna lista. Ideał pierwszy
$\mathfrak p$ pierścienia $\mathbb Z$ zawiera $(6)$
wtedy i tylko wtedy, gdy $6\in\mathfrak p$.
Ideał $(0)$ nie spełnia tego warunku, zaś dla
$\mathfrak p=(p)$, gdzie $p$ jest liczbą pierwszą,
warunek $p\mid6$ pozostawia jedynie $p=2$ lub $p=3$.
Stąd $\mathbb Z/(6)$ ma dokładnie dwa ideały pierwsze:
$(2)/(6)$ i $(3)/(6)$. Oba są maksymalne, ponieważ
ich ilorazy są odpowiednio $\mathbb F_2$ i $\mathbb F_3$.
Są też nieporównywalne, bo $2\notin(3)$ i $3\notin(2)$.
Żaden łańcuch nie ma długości dodatniej, więc
$\dim(\mathbb Z/(6))=0$. Równocześnie $\dim\mathbb Z=1$.
Konkretnie klasy $2$ i $3$ są niezerowymi dzielnikami zera,
gdyż $2\cdot3=0$ modulo $6$.
Przykład pokazuje, że przejście do ilorazu może zmniejszyć
wymiar i że pierścień wymiaru zero może mieć dzielniki zera.
\end{przyklad}



\begin{przyklad}[Wymiar zero nie oznacza, że pierścień jest ciałem]
\label{aa:prz:krull-dual-numbers}
Niech $K$ będzie dowolnym ciałem i
$A=K[T]/(T^2)$, gdzie przez
$\varepsilon=T+(T^2)$ oznaczamy klasę zmiennej $T$.
Każdy element $A$ ma jednoznaczną postać
$a+b\varepsilon$, gdzie $a,b\in K$, oraz
$\varepsilon^2=0$, lecz $\varepsilon\ne0$.
Jeżeli $\mathfrak p$ jest ideałem pierwszym $A$, to
$\varepsilon^2=0\in\mathfrak p$ wymusza
$\varepsilon\in\mathfrak p$. Otrzymujemy
$(\varepsilon)\subseteq\mathfrak p$; z drugiej strony
$A/(\varepsilon)\cong K$ jest ciałem, więc $(\varepsilon)$
jest ideałem maksymalnym i nie może być zawarty w większym
ideale właściwym. Zatem $\mathfrak p=(\varepsilon)$.
Jest to jedyny ideał pierwszy pierścienia $A$, wobec czego
$\dim A=0$. Pierścień $A$ nie jest jednak ciałem:
niezerowy element $\varepsilon$ nie jest odwracalny,
gdyż z $\varepsilon u=1$ wynikałoby
$0=\varepsilon^2u=\varepsilon$, co jest sprzecznością.
Wniosek~\ref{aa:cor:krull-nilradical} wyjaśnia ten wynik
również przez równość $\dim A=\dim(A/(\varepsilon))$.
\end{przyklad}

\begin{przyklad}[Pierścień o nieskończonym wymiarze Krulla]
\label{aa:prz:krull-infinite-dimension}
Niech $K$ będzie ciałem, a
$A=K[X_1,X_2,X_3,\ldots]$ pierścieniem wielomianów
w przeliczalnie wielu zmiennych, przy czym każdy wielomian
zależy od skończonej liczby zmiennych. Dla $r\ge1$
połóżmy $\mathfrak p_r=(X_1,\ldots,X_r)$
oraz $\mathfrak p_0=(0)$. Iloraz
$A/\mathfrak p_r\cong K[X_{r+1},X_{r+2},\ldots]$
nie ma dzielników zera, dlatego każdy $\mathfrak p_r$
jest ideałem pierwszym. Zawieranie
$\mathfrak p_r\subsetneq\mathfrak p_{r+1}$ jest ścisłe:
$X_{r+1}$ nie należy do $\mathfrak p_r$, co widać
z jego niezerowej klasy w tym ilorazie.
Dla każdej liczby naturalnej $m\ge1$ otrzymujemy więc
łańcuch $\mathfrak p_0\subsetneq\cdots\subsetneq
\mathfrak p_m$ długości $m$. Długości łańcuchów
nie są ograniczone, czyli $\dim A=+\infty$.
Przykład nie przeczy późniejszemu obliczeniu wymiaru
$K[X_1,\ldots,X_n]$: tam liczba generatorów jest skończona.
\end{przyklad}
\begin{przyklad}[Lokalizacja może obniżyć wymiar]\label{aa:prz:krull-local-example}
W $\mathbb Z$ wybierzmy $S=\mathbb Z\setminus\{0\}$.
Wtedy $S^{-1}\mathbb Z=\mathbb Q$ ma wymiar $0$,
a $\mathbb Z$ ma wymiar $1$. Z drugiej strony
lokalizacja $\mathbb Z_{(p)}$ względem ideału $(p)$
ma dokładnie dwa ideały pierwsze:
$(0)$ oraz $(p)\mathbb Z_{(p)}$. Wynika to z bijekcji
z ideałami pierwszymi $\mathbb Z$ zawartymi w $(p)$.
Dlatego nadal ma wymiar $1$; odpowiadający im łańcuch
jest ścisły. Nie należy więc oczekiwać, że każda
lokalizacja obniża wymiar: może go zachować.
\end{przyklad}

\begin{uwaga}[Związek z bazami transcendencji]
\label{aa:uw:krull-transcendence-preview}
W pierścieniu $K[X_1,\ldots,X_n]$ występuje łańcuch
$(0)\subsetneq(X_1)\subsetneq\cdots\subsetneq
(X_1,\ldots,X_n)$ długości $n$. To dowodzi nierówności
$\dim K[X_1,\ldots,X_n]\ge n$.
Dla $n=1$ równość $\dim K[X_1]=1$ wynika
z Twierdzenia~\ref{aa:tw:krull-pid}, a dla $n=2$
z Lematu~\ref{aa:lem:krull-plane-dimension}.
Ogólny dowód nierówności odwrotnej dla dowolnego
$n\ge3$ wymaga dodatkowego argumentu, który przeprowadzimy
po rozwinięciu teorii rozszerzeń całkowitych.
Tam też zwiążemy wymiar z
$\operatorname{trdeg}_K\operatorname{Frac}(A)$ dla skończenie
generowanych dziedzin będących $K$-algebrami.
\end{uwaga}
